% Compilador: pdfLaTeX. En Overleaf, selecciona este archivo como documento principal.
% Para incorporar una captura, súbela a la carpeta del .tex con el nombre
% indicado en el recuadro. También puedes cambiar el tercer argumento de \evidencia.
% Compila dos veces si ejecutas pdfLaTeX localmente para actualizar índice y referencias.
\documentclass[12pt,letterpaper]{article}

\usepackage[utf8]{inputenc}
\usepackage[T1]{fontenc}
\usepackage[spanish]{babel}
\usepackage{newtxtext}
\usepackage[letterpaper,margin=1in,headheight=15pt,headsep=0.28in]{geometry}
\usepackage{setspace}
\usepackage{ragged2e}
\usepackage{fancyhdr}
\usepackage{titlesec}
\usepackage{graphicx}
\usepackage{booktabs,tabularx,array}
\usepackage{xurl}
\usepackage[hidelinks]{hyperref}

\doublespacing
\RaggedRight
\setlength{\RaggedRightParindent}{0.5in}
\setlength{\parindent}{0.5in}
\setlength{\parskip}{0pt}
\setlength{\emergencystretch}{2em}
\setcounter{secnumdepth}{0}
\setcounter{tocdepth}{2}

\pagestyle{fancy}
\fancyhf{}
\fancyhead[R]{\thepage}
\renewcommand{\headrulewidth}{0pt}
\fancypagestyle{plain}{%
  \fancyhf{}%
  \fancyhead[R]{\thepage}%
  \renewcommand{\headrulewidth}{0pt}%
}

\titleformat{\section}[block]{\normalfont\normalsize\bfseries\centering}{}{0pt}{}
\titleformat{\subsection}[block]{\normalfont\normalsize\bfseries}{}{0pt}{}
\titlespacing{\section}{0pt}{\baselineskip}{0pt}
\titlespacing{\subsection}{0pt}{\baselineskip}{0pt}

\newcommand{\tituloinforme}{Gestor de incidencias de software con integración continua}
\newcommand{\fechaentrega}{Noviembre de 2026} % Añade el día si lo exige la asignatura.
\newcommand{\entrada}[1]{\par\noindent\hangindent=0.5in\hangafter=1 #1\par}
\newcommand{\notafigura}[1]{\par\noindent\textit{Nota.} #1\par}
\newcommand{\evidencia}[4]{%
  \par\medskip\noindent
  \begin{minipage}{\textwidth}
    \refstepcounter{figure}\label{#1}%
    \noindent\textbf{Figura \thefigure}\par
    \noindent\textit{#2}\par
    \medskip\noindent
    \IfFileExists{#3}{%
      \includegraphics[width=\linewidth,height=3.15in,keepaspectratio]{#3}%
    }{%
      \fbox{\parbox[c][1.45in][c]{\dimexpr\linewidth-2\fboxsep-2\fboxrule\relax}{%
        \centering\small Espacio para la captura\par
        \texttt{#3}\par
      }}%
    }%
    \medskip
    \notafigura{#4}%
  \end{minipage}\par
}
\newcommand{\anexo}[2]{%
  \clearpage
  \section*{Anexo #1\\#2}%
  \addcontentsline{toc}{section}{Anexo #1. #2}%
}

\hypersetup{pdftitle={Gestor de incidencias de software con integración continua},
  pdfauthor={Grupo 5 - Politécnico Grancolombiano}}

\begin{document}

\vspace*{3\baselineskip}
\begin{center}
  \textbf{\tituloinforme}\par
  \textbf{Primera entrega: GitHub y Docker}\par
  \vspace{\baselineskip}
  Lina Bedoya Peña, Jorge Silva Álvarez, Kevin Andrés Molano Gamarra,\\
  Uriel Andrés Rincón y Diego Fernando Gaitán Avella\par
  Facultad de Diseño, Ingeniería e Innovación, Politécnico Grancolombiano\par
  Integración Continua, Grupo 5\par
  Edward Reyes Corredor\par
  \fechaentrega\par
\end{center}

\clearpage
\renewcommand{\contentsname}{Tabla de contenido}
\tableofcontents

\clearpage
\phantomsection
\addcontentsline{toc}{section}{Introducción}
\begin{center}
  \textbf{\tituloinforme}
\end{center}

El equipo desarrolló una aplicación web para registrar incidencias de software, asignarles prioridad y seguir sus cambios de estado. El proyecto sirvió para organizar el trabajo en ramas de Git, integrar una aplicación Flask con PostgreSQL y describir su entorno de ejecución con Docker Compose. El repositorio también incorpora pruebas automáticas que verifican el comportamiento de la aplicación cuando se proponen cambios (Bedoya Peña et al., 2026).

El informe abarca la evolución del prototipo, la implementación integrada en el repositorio y el aporte de documentación y coordinación de Kevin Andrés Molano Gamarra. El estado del repositorio y de las pruebas corresponde al 9 de noviembre de 2026. El Anexo A describe el contrato de datos y el Anexo B reúne los espacios para las evidencias del proceso.

\section{Objetivos}

\subsection{Objetivo general}

Desarrollar un gestor de incidencias con Flask y PostgreSQL y organizar su construcción colaborativa en GitHub, con un entorno reproducible mediante Docker Compose y verificaciones automatizadas mediante GitHub Actions.

\subsection{Objetivos específicos}

\begin{itemize}
  \item Implementar el registro, la consulta, el filtrado, la edición, el cambio de estado y la eliminación de incidencias.
  \item Integrar el esquema de PostgreSQL y comunicar la aplicación y la base de datos mediante Docker Compose.
  \item Automatizar la comprobación de las operaciones del gestor con Pytest y una base de datos dedicada.
  \item Registrar los aportes del equipo con ramas, revisiones y solicitudes de cambio, y documentar el uso del proyecto.
\end{itemize}

\section{Descripción y alcance del proyecto}

El gestor organiza en un solo listado la información que se necesita para dar seguimiento a un error de software. Cada registro contiene un título, una descripción, una prioridad y un estado. El usuario puede elegir una prioridad baja, media o alta. Las incidencias nuevas quedan pendientes y se pueden marcar como en proceso o resueltas. Desde el listado se pueden aplicar filtros, editar los datos y eliminar registros previa confirmación en el navegador.

La entrega cubre la interfaz web, el almacenamiento relacional, la ejecución de dos servicios con Docker Compose y la automatización de pruebas. Se documentan por separado las comprobaciones manuales del prototipo y las verificaciones del sistema integrado.

\section{Marco conceptual}

\subsection{Integración continua y control de versiones}

La integración continua incorpora cambios con frecuencia y ejecuta comprobaciones automáticas para detectar problemas antes de integrar nuevas contribuciones. Git conserva el historial del proyecto y permite trabajar en ramas independientes. En GitHub, un \textit{pull request} reúne los cambios propuestos, la discusión, las revisiones y el estado de las comprobaciones. GitHub Actions permite ejecutar esas comprobaciones a partir de eventos del repositorio (GitHub, s. f.-a, s. f.-b).

\subsection{Aplicación web, persistencia y contenedores}

Flask asocia peticiones HTTP con funciones de Python y permite generar HTML a partir de plantillas (Pallets, s. f.). En el gestor, PostgreSQL almacena las incidencias y Psycopg ejecuta las consultas que necesita Flask. Docker Compose define los servicios y crea una red en la que cada contenedor puede localizar a los demás por el nombre del servicio (Docker, s. f.).

\section{Diseño e implementación}

\subsection{Interfaz y rutas}

El navegador presenta los formularios y el listado generado con Jinja a partir de \texttt{templates/index.html}. Los estilos están en \texttt{static/styles.css}. En \texttt{app.py}, Flask recibe los formularios, valida los campos, ejecuta las consultas y devuelve una respuesta o redirige al listado. La Tabla~\ref{tab:rutas} resume las rutas del sistema integrado. Las Figuras~\ref{fig:prototipo} a~\ref{fig:ui-eliminar} corresponden a la evolución de la interfaz.

\par\medskip\noindent
\begin{minipage}{\textwidth}
  \refstepcounter{table}\label{tab:rutas}
  \noindent\textbf{Tabla \thetable}\par
  \noindent\textit{Rutas HTTP del gestor}\par
  \begin{spacing}{1.15}
  \noindent\begin{tabularx}{\linewidth}{@{} l >{\raggedright\arraybackslash}p{2.48in} >{\raggedright\arraybackslash}X @{}}
    \toprule
    \textbf{Método} & \textbf{Ruta} & \textbf{Operación} \\
    \midrule
    GET & \texttt{/} & Lista o filtra por estado. \\
    POST & \texttt{/incidencias} & Crea una incidencia. \\
    POST & \texttt{/incidencias/<id>/editar} & Edita título, descripción y prioridad. \\
    POST & \texttt{/incidencias/<id>/estado} & Cambia el estado. \\
    POST & \texttt{/incidencias/<id>/eliminar} & Elimina una incidencia. \\
    GET & \texttt{/salud} & Comprueba que Flask responde. \\
    \bottomrule
  \end{tabularx}
  \end{spacing}
  \notafigura{La ruta \texttt{/salud} comprueba el servicio Flask. PostgreSQL tiene su propia comprobación con \texttt{pg\_isready}.}
\end{minipage}\par\medskip

El servidor exige un título y una descripción con contenido, limita el título a 150 caracteres y valida los valores de estado y prioridad. Los datos inválidos producen una respuesta HTTP 400. Si el identificador no existe, la edición, la actualización de estado y la eliminación responden con HTTP 404. Una modificación exitosa redirige al listado con HTTP 302.

\subsection{Persistencia con PostgreSQL}

El prototipo guardaba los registros en una lista de Python y perdía su contenido cuando se reiniciaba Flask. El sistema integrado lee la conexión desde \texttt{DATABASE\_URL}, consulta PostgreSQL con parámetros y conserva los registros en la tabla \texttt{public.incidencias}. El script \texttt{database/gestor\_incidencias.sql} define restricciones para los valores permitidos e índices para estado y prioridad. PostgreSQL asigna el identificador y las fechas al crear cada incidencia. Un disparador actualiza la fecha de modificación cuando cambia un registro.

Por ejemplo, al registrar una incidencia el navegador envía un formulario a \texttt{/incidencias}. Flask valida los datos, ejecuta la inserción con Psycopg y redirige a la página principal. La siguiente petición consulta PostgreSQL y presenta el registro en el listado. El Anexo A especifica los campos que intervienen en ese recorrido.

\subsection{Dockerización y comunicación de servicios}

El archivo \texttt{dockerfile} instala las dependencias y ejecuta Flask sobre una imagen de Python 3.10. En \texttt{docker-compose.yaml}, el servicio \texttt{web} publica el puerto 5000 y el servicio \texttt{db} utiliza PostgreSQL 18 Alpine y publica el 5432. Dentro de la red de Compose, Flask se conecta a \texttt{db:5432}. El servicio web espera a que la comprobación de disponibilidad de PostgreSQL sea satisfactoria.

En la primera inicialización del volumen, la imagen de PostgreSQL ejecuta el script de la base de datos. El volumen \texttt{postgres\_data} conserva los registros entre recreaciones de los contenedores. Para levantar el entorno se copia \texttt{.env.example} a \texttt{.env}, se completan las variables de conexión y se ejecuta \texttt{docker compose up --build -d}. El estado de los servicios se consulta con \texttt{docker compose ps}. Las Figuras~\ref{fig:docker-config} y~\ref{fig:docker-estado} corresponden a la configuración y al estado del entorno. La Figura~\ref{fig:persistencia} reserva un espacio para la comprobación manual de que un registro permanece al reiniciar \texttt{web}.

\section{Pruebas e integración continua}

\subsection{Comprobaciones manuales}

Durante el desarrollo inicial se comprobó desde el navegador el registro, el listado, el filtro por estado, la edición y la eliminación. El registro de Flask mostró redirecciones HTTP 302 para operaciones exitosas seguidas de peticiones GET con respuesta HTTP 200. También se comprobó que crear o editar con un título de 151 caracteres devolvía HTTP 400. Estas comprobaciones corresponden al prototipo de Flask.

\subsection{Pruebas automatizadas y aislamiento}

Pytest comprueba la creación, la edición, el cambio de estado, la eliminación, el listado, los filtros, la persistencia, las fechas y la ruta \texttt{/salud}. La configuración de pruebas exige \texttt{TEST\_DATABASE\_URL} y valida tanto el nombre declarado como el nombre de la base a la que llega la conexión. Solo si el destino es \texttt{gestor\_incidencias\_test} se recrean la tabla y su función de actualización. Las pruebas que usan el cliente Flask vacían la tabla y reinician sus identificadores antes y después de cada caso.

El flujo \texttt{.github/workflows/pruebas.yml} se ejecuta con los \textit{pull requests} y al enviar cambios a \texttt{main}. Utiliza un servicio PostgreSQL 16 y ejecuta una matriz con Python 3.11 y 3.12. Cada trabajo genera un informe JUnit y un archivo XML de cobertura. GitHub Actions registró resultados satisfactorios en ambas versiones tras incorporar el aislamiento (ejecución \href{https://github.com/kmolano-dot/gestor-incidencias/actions/runs/42731745804}{42731745804}) y tras fusionar la actualización del README (ejecución \href{https://github.com/kmolano-dot/gestor-incidencias/actions/runs/42734010046}{42734010046}). La Figura~\ref{fig:acciones} corresponde a los resultados de la matriz.

\section{Organización del equipo y control de versiones}

\subsection{Responsabilidades y aportes}

Uriel Andrés Rincón desarrolló la aplicación Flask y su interfaz. Lina Bedoya Peña definió el esquema de PostgreSQL. Jorge Silva Álvarez preparó Docker Compose y la comunicación de los contenedores. Diego Fernando Gaitán Avella desarrolló y adaptó las pruebas automatizadas y su ejecución en GitHub Actions. Kevin Andrés Molano Gamarra organizó el repositorio, documentó el flujo de colaboración y contribuyó al aislamiento de la base de pruebas.

El historial registra la incorporación de la guía de colaboración y la plantilla de revisión (PR \#1), la aplicación Flask (PR \#2), el esquema SQL (PR \#3), las pruebas iniciales (PR \#4), la conexión entre Flask y PostgreSQL (PR \#5), Docker Compose (PR \#8) y las pruebas sobre PostgreSQL (PR \#9). La rama \texttt{main} también contiene el ajuste de aislamiento del PR \#10 y la actualización del README del PR \#12 (Bedoya Peña et al., 2026). Las Figuras~\ref{fig:repo} a~\ref{fig:pr-flask} corresponden a la vista del repositorio y a los primeros aportes del equipo.

\subsection{Aporte de Kevin Andrés Molano Gamarra}

El aporte inicial consistió en crear el repositorio compartido, invitar a los integrantes y documentar el proceso para trabajar con ramas y revisiones. En el PR \#1 se incorporaron \texttt{CONTRIBUTING.md}, la plantilla de \textit{pull request} y un README con el propósito del proyecto y las responsabilidades. La plantilla pide explicar el cambio y cómo comprobarlo, de modo que el revisor tenga la información necesaria antes de aprobar. El borrador del equipo registra que la regla de protección de \texttt{main} pasó de dos aprobaciones a una. La Figura~\ref{fig:regla} reserva un espacio para mostrar la configuración y la fecha de la captura.

El segundo aporte fue separar la base de datos de la aplicación de la utilizada por Pytest mediante \texttt{TEST\_DATABASE\_URL}. El \href{https://github.com/kmolano-dot/gestor-incidencias/pull/10}{PR \#10}, dirigido a \texttt{test/pruebas-con-postgresql}, incorporó la validación de la base de destino antes de recrear o vaciar los objetos de prueba. El ajuste figura en \texttt{main} en el commit \texttt{ea474d7}. La Figura~\ref{fig:pr-aislamiento} corresponde a esa revisión.

Por último, en la rama \texttt{codex/actualizar-readme}, el commit \texttt{573d6a2} amplió el README con los pasos de arranque en Docker Compose y en el equipo, la configuración de PostgreSQL, las rutas HTTP, el procedimiento de pruebas y los artefactos de GitHub Actions (Molano Gamarra, 2026). El equipo fusionó la actualización en \texttt{main} mediante el \href{https://github.com/kmolano-dot/gestor-incidencias/pull/12}{PR \#12}, en el commit \texttt{d59a9d3}. La Figura~\ref{fig:readme} corresponde a ese cambio de documentación.

\section{Conclusiones}

La aplicación integra una interfaz para gestionar incidencias con un esquema de PostgreSQL que valida sus datos y registra sus fechas. Docker Compose reúne la aplicación y la base de datos en un mismo entorno y define cómo se comunican, mientras que el volumen permite conservar los registros entre recreaciones de los contenedores.

La automatización permite repetir las comprobaciones del gestor con dos versiones de Python. Separar la base de pruebas reduce el riesgo de que la preparación y la limpieza de Pytest afecten los datos de la aplicación. El historial de Git, los \textit{pull requests} y el README permiten rastrear los aportes del equipo y reproducir el entorno.

\clearpage
\section{Referencias}

\entrada{Bedoya Peña, L., Silva Álvarez, J., Molano Gamarra, K. A., Rincón, U. A., \& Gaitán Avella, D. F. (2026). \textit{Gestor de incidencias} (commit d59a9d3) [Código fuente]. GitHub. \url{https://github.com/kmolano-dot/gestor-incidencias/tree/d59a9d3e6886c05bdc856953b2a9820a0ad62f29}}

\entrada{Docker. (s. f.). \textit{Networking in Compose}. Docker Docs. \url{https://docs.docker.com/compose/how-tos/networking/}}

\entrada{GitHub. (s. f.-a). \textit{About pull requests}. GitHub Docs. \url{https://docs.github.com/en/pull-requests/get-started/about-pull-requests}}

\entrada{GitHub. (s. f.-b). \textit{Understanding GitHub Actions}. GitHub Docs. \url{https://docs.github.com/en/actions/get-started/understand-github-actions}}

\entrada{Molano Gamarra, K. A. (2026, 9 de noviembre). \textit{Docs: Actualizar el README con la configuración, las rutas y las pruebas del gestor} [Commit de Git]. GitHub. \url{https://github.com/kmolano-dot/gestor-incidencias/commit/573d6a2}}

\entrada{Pallets. (s. f.). \textit{Quickstart}. Flask Documentation. \url{https://flask.palletsprojects.com/en/stable/quickstart/}}

\anexo{A}{Contrato de integración entre Flask y PostgreSQL}
\setcounter{table}{0}
\renewcommand{\thetable}{A\arabic{table}}
\renewcommand{\theHtable}{A.\arabic{table}}

La tabla \texttt{public.incidencias} reúne los campos que usa Flask y las marcas de tiempo que administra PostgreSQL. Flask envía el título, la descripción y la prioridad al crear un registro. PostgreSQL asigna el identificador, el estado inicial y las fechas. El cambio de estado se procesa en su propia ruta.

\par\medskip\noindent
\begin{minipage}{\textwidth}
  \refstepcounter{table}
  \noindent\textbf{Tabla \thetable}\par
  \noindent\textit{Campos de la tabla de incidencias}\par
  \begin{spacing}{1.15}
  \noindent\begin{tabularx}{\linewidth}{@{} l l >{\raggedright\arraybackslash}X @{}}
    \toprule
    \textbf{Campo} & \textbf{Tipo SQL} & \textbf{Regla} \\
    \midrule
    \texttt{id} & \texttt{BIGINT} & Identidad autogenerada y clave primaria. \\
    \texttt{titulo} & \texttt{VARCHAR(150)} & Obligatorio y con contenido. \\
    \texttt{descripcion} & \texttt{TEXT} & Obligatoria y con contenido. \\
    \texttt{prioridad} & \texttt{VARCHAR(10)} & Baja, Media o Alta. Valor predeterminado en el esquema: Media. \\
    \texttt{estado} & \texttt{VARCHAR(20)} & Pendiente, En proceso o Resuelta. Valor inicial: Pendiente. \\
    \texttt{fecha\_creacion} & \texttt{TIMESTAMPTZ} & Se asigna al insertar el registro. \\
    \texttt{fecha\_actualizacion} & \texttt{TIMESTAMPTZ} & Se asigna al insertar y se actualiza mediante un disparador. \\
    \bottomrule
  \end{tabularx}
  \end{spacing}
  \notafigura{Datos tomados del esquema \texttt{database/gestor\_incidencias.sql} del repositorio del equipo.}
\end{minipage}\par

\anexo{B}{Evidencias gráficas del proceso}
\setcounter{figure}{0}
\renewcommand{\thefigure}{B\arabic{figure}}

\evidencia{fig:prototipo}{Ruta principal del prototipo}{01-prototipo.png}
{La respuesta JSON corresponde a la etapa inicial. La aplicación integrada presenta un listado HTML.}

\evidencia{fig:salud}{Respuesta de la ruta de salud}{02-salud.png}
{La ruta \texttt{/salud} comprueba que Flask está disponible.}

\evidencia{fig:ui-registro}{Formulario y registro de una incidencia}{03-registro.png}
{Formulario de registro e incorporación de una incidencia al listado.}

\evidencia{fig:ui-filtro}{Edición, estados y filtro de incidencias}{04-edicion-filtro.png}
{Edición de los datos de un registro y consulta del listado según el estado seleccionado.}

\evidencia{fig:ui-eliminar}{Confirmación y resultado de eliminar una incidencia}{05-eliminacion.png}
{Confirmación solicitada por la interfaz y estado del listado después de la eliminación.}

\evidencia{fig:docker-config}{Definición de los servicios en Docker Compose}{06-docker-compose.png}
{Muestra la configuración de \texttt{web} y \texttt{db}, el volumen y la comprobación de salud de la base.}

\evidencia{fig:docker-estado}{Servicios en ejecución}{07-docker-estado.png}
{Estado de \texttt{web} y \texttt{db} en Docker Compose.}

\evidencia{fig:persistencia}{Persistencia después de reiniciar Flask}{08-persistencia.png}
{Comparación del mismo identificador antes y después de reiniciar \texttt{web}, con PostgreSQL y el volumen en uso.}

\evidencia{fig:acciones}{Resultados de Pytest en GitHub Actions}{09-github-actions.png}
{Resultados de Pytest para Python 3.11 y 3.12 y artefactos generados por el flujo.}

\evidencia{fig:repo}{Repositorio compartido y rama principal}{10-repositorio.png}
{Muestra la vista de código del repositorio grupal.}

\evidencia{fig:pr-docs}{Pull request de la guía de colaboración}{11-pr-documentacion.png}
{El PR \#1 registra la guía de contribución, la plantilla de revisión y la primera versión del README.}

\evidencia{fig:pr-flask}{Pull request de la aplicación Flask}{12-pr-flask.png}
{El PR \#2 documenta los cambios del componente web y permite mostrar las revisiones solicitadas.}

\evidencia{fig:regla}{Regla de protección de la rama principal}{13-proteccion-main.png}
{Configuración del mínimo de aprobaciones exigidas para integrar cambios en la rama principal.}

\evidencia{fig:pr-aislamiento}{Pull request de aislamiento de las pruebas}{14-pr-aislamiento.png}
{El PR \#10 documenta la separación de la conexión de pruebas y la verificación del nombre de la base.}

\evidencia{fig:readme}{Actualización del README del gestor}{15-readme.png}
{Documentación incorporada en el PR \#12.}

\end{document}
